\documentclass{article}
\usepackage[utf8]{inputenc}
\usepackage[russian]{babel}
\usepackage{amsmath}
\usepackage{amsthm}

\begin{document}

\noindent\textbf{Базовое правило вычислений ($\beta$-редукция):}
\begin{align*}
    (\lambda x. M) N &\rightarrow_\beta M[x \leftarrow N]
\end{align*}

\noindent\textbf{Утверждение.} Данный $\lambda$-терм не имеет нормальной формы.

\begin{proof}[Доказательство]
Для удобства обозначим бесконечно расходящийся субтерм $(\lambda b. b b)(\lambda b. b b)$ как $\Omega$. Выполним последовательную редукцию исходного выражения:
\begin{align*}
    ((\lambda a. (\lambda b. b b) (\lambda b. b b)) b) ((\lambda c. (c b)) (\lambda a. a)) 
    &= ((\lambda b. b b) (\lambda b. b b)) ((\lambda c. (c b)) (\lambda a. a)) \\
    &= \Omega \ ((\lambda c. (c b)) (\lambda a. a)) \\
    &= \Omega \ ((\lambda a. a) b) \\
    &= \Omega \ b
\end{align*}
Так как левый аппликативный хвост сводится к циклическому комбинатору $\Omega$, который бесконечно редуцируется в самого себя, терм расходится и не имеет нормальной формы.
\end{proof}

\end{document}